\section{Аксиомы Пеано}

\begin{axiom}{2.1}{}
$0$ — натуральное число.
\leansource{Nat.instZero}{: Zero Nat}
\end{axiom}

\begin{axiom}{2.2}{}
Если $n$ — натуральное число, то $n{+}{+}$ тоже натуральное число.
\leansource{Nat.succ}{: Nat → Nat}
\end{axiom}

\begin{intuition}
Вместе эти две аксиомы задают порождающий процесс: есть стартовая точка $0$ и правило
«взять следующее», $n \mapsto n{+}{+}$. Сами по себе они ничего не говорят о том, что
получится, если применять правило снова и снова — это уточняют следующие три аксиомы,
исключающие патологии вроде зацикливания (\taoref{ex:2.1.5}{пример 2.1.5}).
\end{intuition}

\begin{remark}
В формализации обе аксиомы выполняются автоматически, без отдельного доказательства.
За это отвечают сами конструкторы индуктивного типа \code{Nat}:
\begin{itemize}
  \item \code{zero : Nat}
  \item \code{succ : Nat → Nat}
\end{itemize}
Тип каждого конструктора сам гарантирует, что результат снова лежит в \code{Nat}.
Следующие три аксиомы устроены иначе. Их придётся доказывать как обычные теоремы.
\end{remark}

\begin{definition}{2.1.3}{}
Определяем $1 := 0{+}{+}$, $2 := 1{+}{+}$, $3 := 2{+}{+}$ и так далее.
\leansource{Nat.instOfNat}{: OfNat Nat n}
\end{definition}

\begin{proposition}{2.1.4}{}
$3$ — натуральное число.
\leansource{Nat.two_succ}{: 2++ = 3}
\end{proposition}

\begin{proof}
$0$ — натуральное число.\by{\taoref{ax:2.1}{аксиома 2.1}}\\
$0{+}{+} = 1$ — натуральное число.\by{\taoref{ax:2.2}{аксиома 2.2}}\\
$1{+}{+} = 2$ — натуральное число.\by{снова \taoref{ax:2.2}{аксиома 2.2}}\\
$2{+}{+} = 3$ — натуральное число.\by{снова \taoref{ax:2.2}{аксиома 2.2}}
\end{proof}

\begin{example}{2.1.5}{}
Рассмотрим систему из четырёх чисел $\{0,1,2,3\}$, в которой операция «следующий»
зацикливается: $0{+}{+}=1$, $1{+}{+}=2$, $2{+}{+}=3$, но $3{+}{+}=0$. Такая система
удовлетворяет \taoref{ax:2.1}{аксиоме 2.1} и \taoref{ax:2.2}{аксиоме 2.2}, но не
соответствует привычному представлению о натуральных числах. Похожее происходит и на
практике: если целочисленный счётчик в компьютере переполняется, он резко возвращается
к нулю.
\end{example}

\begin{axiom}{2.3}{}
$0$ не является последователем никакого натурального числа: для любого натурального
$n$ выполняется $n{+}{+} \neq 0$.
\leansource{Nat.succ_ne}{(n : Nat) : n++ ≠ 0}
\end{axiom}

\begin{intuition}
Без этой аксиомы цепочка $0,\,0{+}{+},\,(0{+}{+}){+}{+},\,\ldots$ могла бы замкнуться сама
на себя, как в \taoref{ex:2.1.5}{примере 2.1.5} с четырёхэлементным циклом. Требуя, чтобы
$0$ никогда не оказалось чьим-то последователем, аксиома запрещает именно этот вид петли —
$0$ обязано быть самой первой точкой цепочки.
\end{intuition}

\begin{remark}
Эта аксиома как раз запрещает зацикливание из \taoref{ex:2.1.5}{примера 2.1.5}.
В Lean она доказывается. Конструкторы индуктивного типа устроены так, что \code{zero}
никогда не совпадает с результатом \code{succ}, и тактика \code{injection} превращает это
различие в противоречие.
\end{remark}

\begin{proposition}{2.1.6}{}
$4 \neq 0$.
\leansource{Nat.four_ne}{: (4 : Nat) ≠ 0}
\end{proposition}

\begin{proof}
$4 = 3{+}{+}$.\by{\taoref{def:2.1.3}{определение 2.1.3}}\\
$3$ — натуральное число.\by{\taoref{ax:2.1}{аксиома 2.1}, \taoref{ax:2.2}{аксиома 2.2}}\\
$3{+}{+} \neq 0$, то есть $4 \neq 0$.\by{\taoref{ax:2.3}{аксиома 2.3}}
\end{proof}

\begin{remark}
Из одного лишь определения числа $4$ это не следует — понадобилась
\taoref{ax:2.3}{аксиома 2.3}. В патологической системе из
\taoref{ex:2.1.5}{примера 2.1.5} утверждение вообще ложно: там $4 = 0$.
\end{remark}

\begin{example}{2.1.7}{}
\taoref{ax:2.1}{Аксиом 2.1}–\taoref{ax:2.3}{2.3} всё ещё недостаточно. Рассмотрим систему из пяти чисел $\{0,1,2,3,4\}$,
в которой операция «следующий» упирается в потолок на $4$: $0{+}{+}=1$, \dots, $3{+}{+}=4$,
но $4{+}{+}=4$ (то есть $5=4$, $6=4$ и так далее) — это не противоречит ни одной из трёх
аксиом. Похожая проблема возникает и при зацикливании в произвольное другое число,
например $4{+}{+}=1$.
\end{example}

\begin{remark}
За этим списком аксиом стоит общий метод построения аксиоматической системы. Начинают с
минимального набора аксиом, а затем рассматривают конкретные модели, которые ему
удовлетворяют. Если среди таких моделей находится что-то, противоречащее интуитивным
представлениям о натуральных числах — как зацикливание в \taoref{ex:2.1.5}{примере 2.1.5}
или упирание в потолок в \taoref{ex:2.1.7}{примере 2.1.7}, — в набор добавляется ровно та
аксиома, которая нужна, чтобы исключить именно эту патологию, и не более того.
\end{remark}

\begin{axiom}{2.4}{}
Разные натуральные числа имеют разных последователей: если $n \neq m$, то
$n{+}{+} \neq m{+}{+}$. Равносильно: если $n{+}{+} = m{+}{+}$, то $n = m$.
\leansource{Nat.succ_ne_succ}{(n m : Nat) : n ≠ m → n++ ≠ m++}
\leansource{Nat.succ_cancel}{(hnm : n++ = m++) : n = m}
\end{axiom}

\begin{intuition}
\taoref{ax:2.3}{Аксиома 2.3} запрещает цепочке замыкаться на $0$, но не мешает ей склеиться где-то дальше или
упереться в потолок, как в \taoref{ex:2.1.7}{примере 2.1.7} ($4{+}{+}=4$ или $4{+}{+}=1$).
Инъективность последователя сама по себе такую склейку впрямую не запрещает — она лишь даёт
способ её \emph{размотать}: увидев $n{+}{+}=m{+}{+}$, всегда можно снять слой «++» сразу с
обеих сторон и перейти к $n=m$. Применённая к $4{+}{+}=4$, эта операция шаг за шагом сводит
равенство к $4=3$, затем к $3=2$, $2=1$ и наконец к $1=0$ — а это в точности $0{+}{+}=0$,
что запрещено \taoref{ax:2.3}{аксиомой 2.3}. Так же $4{+}{+}=1$ сразу разматывается до $3=0$ —
опять недопустимо. Всю настоящую работу здесь делает \taoref{ax:2.3}{аксиома 2.3}, \taoref{ax:2.4}{аксиома 2.4}
лишь даёт инструмент, чтобы до неё добраться.
\end{intuition}

\begin{remark}
Эта аксиома исключает обе патологии из \taoref{ex:2.1.7}{примера 2.1.7}: теперь
$4{+}{+} \neq 4$ и $4{+}{+} \neq 1$.
\end{remark}

\begin{proposition}{2.1.8}{}
$6 \neq 2$.
\leansource{Nat.six_ne_two}{: (6 : Nat) ≠ 2}
\end{proposition}

\begin{proof}
Предположим противное: $6 = 2$.\\
Тогда $5{+}{+} = 1{+}{+}$, откуда $5 = 1$.\by{\taoref{ax:2.4}{аксиома 2.4}}\\
Тогда $4{+}{+} = 0{+}{+}$, откуда $4 = 0$.\by{\taoref{ax:2.4}{аксиома 2.4}}\\
Это противоречит \taoref{prop:2.1.6}{утверждению 2.1.6}.
\end{proof}

\begin{axiom}{2.5}{принцип математической индукции}
Пусть $P(n)$ — свойство, применимое к натуральному числу $n$. Если $P(0)$ истинно и для
каждого $n$ из истинности $P(n)$ следует истинность $P(n{+}{+})$, то $P(n)$ истинно для
всех натуральных $n$.
\leansource{Nat.induction}{(P : Nat → Prop) (hbase : P 0) (hind : ∀ n, P n → P (n++))\leanbreak: ∀ n, P n}
\end{axiom}

\begin{intuition}
\taoref{ax:2.3}{Аксиомы 2.3} и \taoref{ax:2.4}{2.4} гарантируют, что цепочка $0,0{+}{+},(0{+}{+}){+}{+},\ldots$ бесконечна и
не повторяется — но сами по себе не запрещают числовой системе содержать что-то ещё в
стороне от этой цепочки, недостижимое из $0$ конечным числом шагов «взять следующее». Именно
индукция закрывает эту лазейку: свойство, верное в $0$ и переходящее через каждый шаг,
обязано покрыть всю систему целиком, а значит и «лишних» элементов вне цепочки быть не может.
\end{intuition}

\begin{remark}
Строго говоря, это схема аксиом: для каждого свойства $P$ — своя аксиома. Она устроена
иначе, чем \taoref{ax:2.1}{аксиомы 2.1}–\taoref{ax:2.4}{2.4}: те говорят о конкретных
числах, эта — о свойствах. Именно она исключает «лишние» элементы вроде $0.5$:
без неё ничто не мешает числовой системе содержать что-то помимо $0$ и чисел, полученных из
него прибавлением единицы. В Lean её заменяет тактика \code{induction}.
\end{remark}

\begin{assumption}{2.6}{существование натуральных чисел}
Предполагается, что существует числовая система $\mathbb{N}$, элементы которой называются
натуральными числами и для которой выполняются \taoref{ax:2.1}{аксиомы 2.1}–\taoref{ax:2.5}{2.5}.
\leansource{Nat}{zero : Nat, succ : Nat → Nat}
\end{assumption}

\begin{remark}
$\mathbb{N}$ здесь строится явно — индуктивным типом \code{Nat}. Чисто аксиоматический
подход, через теорию множеств, появится в эпилоге этой главы. Оба подхода задают одну и ту
же структуру с точностью до изоморфизма — какую бы модель ни выбрать, лишь бы она
подчинялась аксиомам, для математики этого достаточно.
\end{remark}

\begin{proposition}{2.1.16}{рекурсивные определения}
Пусть для каждого натурального $n$ задана функция $f_n : \mathbb{N} \to \mathbb{N}$, и пусть
$c \in \mathbb{N}$. Тогда существует единственная последовательность $(a_n)_{n \in \mathbb{N}}$
такая, что $a_0 = c$ и $a_{n{+}{+}} = f_n(a_n)$ для всех $n$.
\leansource{Nat.recurse_uniq}{(f : Nat → Nat → Nat) (c : Nat)\leanbreak: ∃! (a : Nat → Nat), a 0 = c ∧ ∀ n, a (n++) = f n (a n)}
\end{proposition}

\begin{intuition}
\taoref{ax:2.1}{Аксиомы 2.1}–\taoref{ax:2.5}{2.5} сами по себе не дают права писать рекурсивные определения вроде «$a_0=c$,
$a_{n{+}{+}}=f_n(a_n)$» — они говорят только про $0$ и «следующее», ничего впрямую не
утверждая про произвольные последовательности. Это утверждение и есть недостающее звено: оно доказывает, что такая
последовательность существует и единственна, опираясь именно на индукцию (\taoref{ax:2.5}{аксиома 2.5}) —
дальше на него можно ссылаться каждый раз, когда определение вводится «по индукции», как,
например, сложение в следующем разделе.
\end{intuition}

\begin{remark}
За существование отвечает явная конструкция \code{Nat.recurse}: рекурсивная функция, которая
сопоставляет $0$ значение $c$, а $n{+}{+}$ — значение $f_n$ от уже построенного $a_n$. То, что
она действительно удовлетворяет обоим условиям на каждом шаге, проверяется одной редукцией
определения (\code{Nat.recurse\_zero}, \code{Nat.recurse\_succ}). Единственность требует
настоящего применения \taoref{ax:2.5}{аксиомы 2.5}: если $a$ и \code{Nat.recurse f c}
совпадают в $0$ и одинаково устроен переход к следующему элементу, они совпадают всюду.
\end{remark}

\begin{remark}
Определение натуральных чисел здесь аксиоматическое: мы не говорим, что такое натуральные
числа, и не отвечаем на вопросы вроде «из чего они сделаны», «физические ли это объекты» или
«что они измеряют». Вместо этого перечислены некоторые действия, которые с ними можно
делать — на самом деле единственная операция, которую мы пока определили, это увеличение на
единицу, — и некоторые их свойства. Так работает математика: она рассматривает свои объекты
абстрактно, её интересуют только свойства, которыми эти объекты обладают. Неважно, означает
ли натуральное число определённое расположение косточек на счётах, определённую организацию
битов в памяти компьютера или какой-то более абстрактный концепт без физического носителя.
Достаточно, чтобы их можно было увеличивать, проверять на равенство, а позже выполнять с ними
сложение и умножение — тогда они годятся в качестве чисел для математических целей,
разумеется, при условии, что подчиняются нужным аксиомам. Натуральные числа можно
сконструировать из других математических объектов, например из множеств, но способов
построить рабочую модель натуральных чисел несколько. Спорить о том, какая модель «истинная»,
с точки зрения математика бессмысленно: достаточно, чтобы модель удовлетворяла всем аксиомам
и работала как нужно.
\end{remark}

\begin{remark}
Исторически осознание того, что числа можно изучать аксиоматически, появилось совсем
недавно — чуть больше ста лет назад. Раньше числа обычно жёстко привязывали к чему-то
внешнему: к количеству элементов в множестве, к длине отрезка, к массе физического объекта и
так далее. Это более-менее работало, пока не приходилось переходить от одной числовой системы
к другой. Например, представлять числа как счётные шарики отлично подходит для чисел $3$ и
$5$, но плохо годится для $-3$, $1/3$, $\sqrt2$ или $3{+}4i$. Поэтому каждый крупный шаг в
теории чисел — отрицательные числа, иррациональные числа, комплексные числа, даже ноль —
вызывал много ненужных философских терзаний. Великое открытие конца XIX века в том, что числа
можно понимать абстрактно, через аксиомы, не требуя обязательной конкретной модели. Конечно,
математик может использовать любую из таких моделей, когда это удобно для интуиции и
понимания, но их же можно без сожаления отбросить, если они начинают мешать.
\end{remark}
